%! Algebra 1 Priority Standards Quiz (KEY)
% GENERATED BODY: the blank and key bodies are byte-identical except the header title;
% answers live in \slot / \optok / work, which the preamble shows or hides.
% Standards: 2023 VA SOL Algebra 1 — A.EO.1, A.EO.2, A.EO.3, A.EO.4, A.EI.1, A.EI.2, A.EI.3,
% A.F.1, A.F.2. A.ST.1 (regression, scatterplots) is deliberately left out.
\documentclass[12pt]{article}
\usepackage{algebra2-article}
\usepackage{algebra2-key}   % pulls in -boxes; do NOT also load -boxes
% Coordinate grid: {xmin}{xmax}{ymin}{ymax}
\newcommand{\lgrid}[4]{%
  \draw[step=1, linegray!40, very thin] (#1,#3) grid (#2,#4);
  \draw[->, semithick, charcoal] (#1,0)--(#2,0) node[below right, font=\scriptsize]{$x$};
  \draw[->, semithick, charcoal] (0,#3)--(0,#4) node[left, font=\scriptsize]{$y$};}
% Part-divider strip (headlinebox takes one arg: the background color)
\newcommand{\parthead}[1]{%
  \boxguard[10]\vspace{6pt}\noindent
  \begin{headlinebox}{forest}\color{white}\bfseries #1\end{headlinebox}%
  \vspace{4pt}}
% THE WORK RULE for a test: \slot{width}{answer} and \opt / \optok are the ONLY
% lines that differ between blank and key, and they differ in the preamble, not the body.
\newcommand{\slot}[2]{\underline{\makebox[#1]{\ans{#2}}}}
\newcommand{\opt}[2]{(#1)~#2}
\newcommand{\optok}[2]{\mbox{\llap{\textcolor{keyred}{$\boldsymbol{\rightarrow}$}\,}\textcolor{keyred}{\textbf{(#1)}~#2}}}

\begin{document}

\pageheader{Algebra 1 Review}{Priority Standards Quiz --- Answer Key}
\namedateperiod

\vspace{0.06in}
\noindent\textbf{Instructions:} No calculator. Show your work where space is given. Circle the
letter for multiple choice. Leave radicals in simplest form.
\hfill\textbf{25 questions, 4 pts each = 100 pts}

% ======================================================================
\parthead{Part A --- Expressions, Exponents \& Radicals \hfill\normalfont\small A.EO.1, A.EO.3, A.EO.4}

\begin{enumerate}[leftmargin=1.9em, itemsep=10pt, label=\arabic*.]
  \item Evaluate $2x^2-|x-5|$ when $x=-3$.
    \begin{work}
      2(-3)^2-|-3-5| &= 2(9)-|-8| \\
      &= 18-8
    \end{work}
    \hfill Value: \slot{2.4cm}{$10$}

  \item Which expression represents ``five less than three times a number $n$''?
    \\[4pt]
    \opt{A}{$5-3n$} \hfill \optok{B}{$3n-5$} \hfill \opt{C}{$3(n-5)$} \hfill \opt{D}{$5n-3$}

  \item Simplify $(2x^3y)^2\cdot x^{-4}$. Use only positive exponents.
    \begin{work}
      (2x^3y)^2\cdot x^{-4} &= 4x^6y^2\cdot x^{-4}
    \end{work}
    \hfill Answer: \slot{3.0cm}{$4x^2y^2$}

  \item Simplify $\dfrac{12a^5b^2}{4a^2b^6}$. Use only positive exponents.
    \begin{work}
      \frac{12}{4}\cdot a^{5-2}\cdot b^{2-6} &= 3a^3b^{-4}
    \end{work}
    \hfill Answer: \slot{3.0cm}{$\dfrac{3a^3}{b^4}$}

  \item Simplify $\sqrt{72}$.
    \begin{work}
      \sqrt{72} &= \sqrt{36}\cdot\sqrt{2}
    \end{work}
    \hfill Answer: \slot{3.0cm}{$6\sqrt{2}$}

  \item Which is equivalent to $2\sqrt{3}+5\sqrt{12}$?
    \\[4pt]
    \opt{A}{$7\sqrt{15}$} \hfill \opt{B}{$7\sqrt{3}$} \hfill \optok{C}{$12\sqrt{3}$} \hfill
    \opt{D}{$22\sqrt{3}$}
\end{enumerate}

% ======================================================================
\parthead{Part B --- Polynomials \& Factoring \hfill\normalfont\small A.EO.2}

\begin{enumerate}[leftmargin=1.9em, itemsep=10pt, label=\arabic*., start=7]
  \item Subtract: $(3x^2-5x+2)-(x^2+2x-6)$.
    \begin{work}
      &= 3x^2-5x+2-x^2-2x+6
    \end{work}
    \hfill Answer: \slot{3.6cm}{$2x^2-7x+8$}

  \item Multiply: $(2x-3)(x+4)$.
    \begin{work}
      &= 2x^2+8x-3x-12
    \end{work}
    \hfill Answer: \slot{3.6cm}{$2x^2+5x-12$}

  \item Factor: $x^2-5x-14$.
    \\[1.4cm]
    \hspace*{\fill}Answer: \slot{3.6cm}{$(x-7)(x+2)$}

  \item Factor: $2x^2+7x+3$.
    \\[1.4cm]
    \hspace*{\fill}Answer: \slot{3.6cm}{$(2x+1)(x+3)$}

  \item Which is the \textbf{complete} factorization of $3x^2-27$?
    \\[4pt]
    \opt{A}{$3(x^2-9)$} \hfill \opt{B}{$3(x-3)^2$} \hfill \optok{C}{$3(x-3)(x+3)$} \hfill
    \opt{D}{$(3x-9)(x+3)$}
\end{enumerate}

% ======================================================================
\parthead{Part C --- Linear Equations, Inequalities \& Systems \hfill\normalfont\small A.EI.1, A.EI.2}

\begin{enumerate}[leftmargin=1.9em, itemsep=10pt, label=\arabic*., start=12]
  \item Solve: $4(x-2)=2x+6$.
    \begin{work}
      4x-8 &= 2x+6 \\
      2x &= 14
    \end{work}
    \hfill Solution: \slot{2.4cm}{$x=7$}

  \item Which is the solution of $-3x+5\le 17$?
    \\[4pt]
    \opt{A}{$x\le -4$} \hfill \optok{B}{$x\ge -4$} \hfill \opt{C}{$x\ge 4$} \hfill
    \opt{D}{$x\le 4$}

  \item Solve $A=\tfrac{1}{2}bh$ for $h$.
    \begin{work}
      2A &= bh
    \end{work}
    \hfill $h=$ \slot{2.4cm}{$\dfrac{2A}{b}$}

  \item How many solutions does $2(x+3)=2x+5$ have?
    \\[4pt]
    \opt{A}{exactly one} \hfill \optok{B}{no solution} \hfill \opt{C}{exactly two} \hfill
    \opt{D}{infinitely many}

  \item Solve the system. \quad
    $\begin{cases} y=2x-1 \\ x+y=8 \end{cases}$
    \begin{work}
      x+(2x-1) &= 8 \\
      3x &= 9 \quad\Rightarrow\quad x=3 \\
      y &= 2(3)-1 = 5
    \end{work}
    \hfill Solution: \slot{2.8cm}{$(3,\,5)$}

  \item Which ordered pair is a solution of $y>-x+2$?
    \\[4pt]
    \opt{A}{$(0,\,0)$} \hfill \opt{B}{$(1,\,1)$} \hfill \optok{C}{$(3,\,0)$} \hfill
    \opt{D}{$(-1,\,2)$}
\end{enumerate}

% ======================================================================
\parthead{Part D --- Linear Functions \hfill\normalfont\small A.F.1}

\begin{enumerate}[leftmargin=1.9em, itemsep=10pt, label=\arabic*., start=18]
  \item For the line $3x-2y=8$, find the slope and the $y$-intercept.
    \begin{work}
      -2y &= -3x+8 \\
      y &= \frac{3}{2}x-4
    \end{work}
    \hfill Slope: \slot{2.0cm}{$\tfrac{3}{2}$} \qquad $y$-intercept: \slot{2.0cm}{$(0,\,-4)$}

  \item Write the equation, in $y=mx+b$ form, of the line \textbf{parallel} to $y=-2x+5$ that
    passes through $(1,\,4)$.
    \begin{work}
      y-4 &= -2(x-1) \\
      y-4 &= -2x+2
    \end{work}
    \hfill Equation: \slot{3.6cm}{$y=-2x+6$}

  \boxguard[9]
  \item Which equation matches the graph?
    \\[4pt]
    \begin{minipage}[c]{0.36\linewidth}\centering
    \begin{tikzpicture}[scale=0.36]
      \lgrid{-4}{4}{-4}{4}
      \draw[very thick, navy] (-0.5,4)--(3.5,-4);
      \fill[navy] (0,3) circle (4pt); \fill[navy] (2,-1) circle (4pt);
    \end{tikzpicture}
    \end{minipage}%
    \begin{minipage}[c]{0.6\linewidth}
    \opt{A}{$y=2x+3$} \\[6pt]
    \opt{B}{$y=-\tfrac{1}{2}x+3$} \\[6pt]
    \optok{C}{$y=-2x+3$} \\[6pt]
    \opt{D}{$y=-2x-3$}
    \end{minipage}

  \item For $f(x)=-3x+12$, find $f(-2)$, and find the value of $x$ for which $f(x)=0$.
    \begin{work}
      f(-2) &= -3(-2)+12 = 18 \\
      0 &= -3x+12 \quad\Rightarrow\quad x=4
    \end{work}
    \hfill $f(-2)=$ \slot{1.6cm}{$18$} \qquad $x=$ \slot{1.6cm}{$4$}
\end{enumerate}

% ======================================================================
\parthead{Part E --- Functions \& Quadratics \hfill\normalfont\small A.F.2, A.EI.3}

\begin{enumerate}[leftmargin=1.9em, itemsep=10pt, label=\arabic*., start=22]
  \item Which relation is \textbf{not} a function?
    \\[4pt]
    \opt{A}{$\{(1,2),\,(2,2),\,(3,5)\}$} \hfill \optok{B}{$\{(2,1),\,(2,4),\,(3,5)\}$}
    \\[6pt]
    \opt{C}{$\{(0,0),\,(1,1),\,(2,4)\}$} \hfill \opt{D}{$\{(-1,3),\,(0,3),\,(1,3)\}$}

  \boxguard[12]
  \item The graph of $f(x)=x^2-2x-3$ is shown. Read its key features.
    \\[4pt]
    \begin{minipage}[c]{0.40\linewidth}\centering
    \begin{tikzpicture}[scale=0.36]
      \lgrid{-4}{5}{-5}{5}
      \draw[very thick, navy, domain=-2:4, smooth, samples=40] plot (\x, {\x*\x-2*\x-3});
      \foreach \p in {(-1,0),(3,0),(1,-4),(0,-3)} \fill[navy] \p circle (4pt);
    \end{tikzpicture}
    \end{minipage}%
    \begin{minipage}[c]{0.58\linewidth}
    Zeros: \slot{3.4cm}{$x=-1$, \ $x=3$} \\[12pt]
    Vertex: \slot{2.4cm}{$(1,\,-4)$} \\[12pt]
    $y$-intercept: \slot{2.4cm}{$(0,\,-3)$} \\[12pt]
    Range: \slot{2.4cm}{$y\ge -4$}
    \end{minipage}

  \item Solve by factoring: $x^2+2x-15=0$.
    \begin{work}
      (x+5)(x-3) &= 0 \\
      x+5=0 \quad&\text{or}\quad x-3=0
    \end{work}
    \hfill Solutions: \slot{3.6cm}{$x=-5$ or $x=3$}

  \item Solve: $x^2-4x-1=0$. Give exact answers in simplest radical form.
    \begin{work}
      x &= \frac{4\pm\sqrt{(-4)^2-4(1)(-1)}}{2(1)} = \frac{4\pm\sqrt{20}}{2} \\
      &= \frac{4\pm 2\sqrt{5}}{2}
    \end{work}
    \hfill Solutions: \slot{3.6cm}{$x=2\pm\sqrt{5}$}
\end{enumerate}

\end{document}
